% SYNTHETIC EVALUATION DRAFT — written for the retrieval eval in evals/.
% Not a real paper; contains no experimental results.
\documentclass{article}
\title{Risk-Constrained Deep Reinforcement Learning for Portfolio Allocation and Order Execution}
\begin{document}
\maketitle

\begin{abstract}
We study deep reinforcement learning agents for two financial decision
problems: allocating capital across a basket of equities and executing large
orders in a limit order book with minimal market impact. Standard reward
formulations that maximize cumulative return lead to policies with excessive
drawdowns and turnover. This draft formulates trading as a constrained Markov
decision process with a conditional value-at-risk constraint and transaction
costs, and proposes a Lagrangian actor-critic method trained in a market
simulator calibrated to historical data.
\end{abstract}

\section{Introduction}
Algorithmic trading systems must make sequential decisions under uncertainty:
how much of each asset to hold, when to rebalance, and how to split a large
parent order into child orders over time. Reinforcement learning (RL) is a
natural fit because the agent's actions influence future states through
transaction costs, inventory and price impact.

However, applying deep RL to trading raises specific difficulties. Financial
returns are non-stationary and have a low signal-to-noise ratio, so agents
overfit easily to the backtest period. Rewards based purely on profit and loss
ignore risk, and a policy that is profitable on average may suffer
unacceptable drawdowns. Finally, realistic execution requires modelling the
limit order book, slippage and the agent's own market impact.

Our contributions are: (1) a constrained MDP formulation of portfolio
allocation with a CVaR constraint on episode returns and proportional
transaction costs; (2) a Lagrangian proximal policy optimization (PPO) agent
with a distributional critic that estimates the return quantiles needed for
the CVaR constraint; and (3) a walk-forward backtesting protocol designed to
reduce look-ahead bias.

\section{Related Work}
Early work applied recurrent reinforcement learning to optimize the Sharpe
ratio directly. Deep Q-networks and policy gradient methods have since been
used for single-asset trading, and actor-critic methods such as DDPG, PPO and
SAC for continuous portfolio weights. Optimal execution has a long history in
the Almgren--Chriss framework, and RL approaches learn execution schedules from
limit order book features. Risk-sensitive RL optimizes objectives such as
mean-variance or CVaR, and constrained RL methods use Lagrangian relaxation to
enforce cost constraints. Distributional RL learns the full return
distribution rather than its mean.

\section{Problem Formulation}
\subsection{State, Action and Reward}
The state contains a window of past log returns, technical indicators, current
portfolio weights and cash. The action is a vector of target portfolio weights
on the probability simplex. The reward at each step is the log portfolio
return net of proportional transaction costs on the traded volume.

\subsection{Risk Constraint}
We require the conditional value-at-risk at level $\alpha$ of the episode
return to stay above a threshold. The constrained objective is relaxed with a
Lagrange multiplier that is updated by dual gradient ascent.

\section{Method}
The policy and critic share a temporal convolutional encoder over the return
window. The critic is a quantile network that outputs a set of return
quantiles, from which the CVaR of the current policy is estimated. The actor is
trained with the PPO clipped objective on the Lagrangian reward. For order
execution, the same agent architecture acts on limit order book snapshots and
chooses the size and price level of child orders, with a penalty for unfilled
inventory at the end of the horizon.

\section{Experimental Plan}
We plan to train in a market simulator replaying historical daily equity data
and a limit order book simulator with an agent-based price impact model.
Evaluation will use walk-forward splits in which the agent is retrained on an
expanding window and tested on the following period. Baselines include
equal-weight and mean-variance portfolios, unconstrained PPO and SAC, and the
Almgren--Chriss execution schedule. Metrics are annualized return, Sharpe
ratio, maximum drawdown, turnover and implementation shortfall.

\section{Limitations}
Simulator calibration determines how well policies transfer to live markets,
and agent-based impact models may not capture rare liquidity events. The
backtest cannot account for the agent's effect on other market participants.

\end{document}
